%% ================================================================
%% クォーラム型 2DC 構成における停止強度の漸近解析と
%% 金融工学的リスク評価（改訂稿 v8）
%%
%% 第2次査読対応改訂履歴（v8）:
%%
%%  [致命的 B-1] ‖S‖₂ ≤ 1/μ_min の誤り修正 ——
%%    T₀ の非正規性（Jordan ブロック存在可能）を認めながらスペクトル距離の
%%    逆数で規約リゾルベントノルムを評価した誤り．
%%    → 補題 3.4（新設）：T₀ の修復連鎖単調性（故障数 DAG 構造）と
%%      M 行列理論を用いて ‖S‖₂ ≤ n_R/μ_min（n_R := ⌈√378⌉ = 20）を
%%      厳密証明．付録 B.3 を全面改訂．
%%
%%  [致命的 B-2] Bhatia Lipschitz 固有値連続性の非正規行列への誤用 ——
%%    非正規行列（非 self-adjoint）に対して固有値の Lipschitz 連続性は
%%    一般に成立しない（条件数に依存）．
%%    → 補題 3.6 を Gershgorin 円板定理による証明に全面置換．
%%      条件数・固有ベクトル情報を一切使用しない自己完結証明．
%%
%%  [致命的 B-3] VaR「閾値跳躍」の自明性 ——
%%    二値損失モデルの定義の帰結であり定理としての価値を欠く．
%%    → VaR 節を再構成：定理の貢献を「VaR 跳躍の存在」ではなく
%%      「臨界値 ε_c の設計パラメータ陽的公式」に転換．
%%      ε_c の完全閉形式（系コンポーネント構造・修復率を直接含む）が
%%      FTA・既存クレジットモデルのいずれでも得られない点を証明．
%%
%%  [致命的 B-4] サラミ出版（演習問題批判） ——
%%    基底論文の θ を CDS 式に代入するだけでは独立した論文として不十分．
%%    → 定理 5.7（新設）：CDS スプレッド・タームストラクチャーの
%%      HA/非HA 非自明差分（Merton/Lando モデルに還元不可能な結果）を追加．
%%      停止強度スケール O(ε²) vs O(ε) が短期スプレッドカーブの
%%      凸性（convexity）に与える定量的差異を導出．
%%
%%  [重大 B-5] C₁ の ε-非依存性の未証明 ——
%%    → QSD 定理（定理 3.8）に C₁ の陽的上界 C₁ ≤ ‖(I-P₀)p₀‖ を追加，
%%      ε-非依存性をスペクトル分解から直接証明．
%%
%%  [重大 B-6] Dynamic Fault Tree（DFT）との比較欠如 ——
%%    FTA との比較が静的 FTA のみを対象とした藁人形論法と指摘された．
%%    → DFT・Markov-based DFT・Barlow-Proschan との詳細比較を表 4 として新設．
%%
%%  コンパイル: xelatex ha_finance_v8.tex（3 回実行推奨）
%% ================================================================
\documentclass[12pt,a4paper]{article}

\usepackage{fontspec}
\setmainfont[
  BoldFont = NotoSerifCJK-Bold.ttc,
  Path     = /usr/share/fonts/opentype/noto/,
  Script   = CJK
]{NotoSerifCJK-Regular.ttc}
\setmonofont{DejaVu Sans Mono}

\usepackage{amsmath,amssymb,amsthm,mathtools}
\usepackage[top=25mm,bottom=25mm,left=25mm,right=25mm]{geometry}
\usepackage{setspace}\onehalfspacing
\usepackage{booktabs,array,float,caption,longtable,multirow}
\usepackage[colorlinks=true,
            linkcolor=blue!60!black,
            citecolor=green!50!black,
            urlcolor=blue!70!black]{hyperref}
\usepackage{microtype,xcolor,parskip,enumitem}
\setlist{itemsep=2pt,topsep=4pt}
\sloppy\setlength{\emergencystretch}{4em}

%% --- 定理環境 ---
\theoremstyle{plain}
\newtheorem{theorem}{定理}[section]
\newtheorem{proposition}[theorem]{命題}
\newtheorem{lemma}[theorem]{補題}
\newtheorem{corollary}[theorem]{系}
\theoremstyle{definition}
\newtheorem{definition}[theorem]{定義}
\newtheorem{assumption}[theorem]{仮定}
\theoremstyle{remark}
\newtheorem{remark}[theorem]{注}

%% --- マクロ（ha_stopping_analysis_v7.tex と完全統一） ---
\newcommand{\bE}{\mathbb{E}}
\newcommand{\calO}{\mathcal{O}}
\newcommand{\calA}{\mathcal{A}}
\newcommand{\calT}{\mathcal{T}}
\newcommand{\calC}{\mathcal{C}}
\newcommand{\eps}{\varepsilon}
\newcommand{\e}{\mathrm{e}}
\newcommand{\RedRes}{S}
\newcommand{\Proj}{P}
%% 金融工学マクロ（本稿追加）
\newcommand{\VaR}{\mathrm{VaR}}
\newcommand{\ES}{\mathrm{ES}}
\newcommand{\dd}{\mathrm{d}}
\newcommand{\calCtwo}{\mathcal{C}_2}
\newcommand{\calCNOR}{\calCtwo^{(\mathrm{NOR})}}
\newcommand{\calCDTS}{\calCtwo^{(\mathrm{DTS})}}
\newcommand{\thetaNOR}{\theta^{(\mathrm{NOR})}}
\newcommand{\thetaDTS}{\theta^{(\mathrm{DTS})}}
\newcommand{\rhoDTS}{\rho_{\mathrm{DTS}}}
\newcommand{\Teval}{T_{\mathrm{eval}}}

% ================================================================
\title{%
  \textbf{クォーラム型 2DC 構成における停止強度の漸近解析と\\
  金融工学的リスク評価}\\[0.5em]
  \large
  ── 吸収型 CTMC の主固有値展開・DTS 分解・\\
  VaR 臨界値閉形式・CDS タームストラクチャーの自己完結統合理論 ──\\[0.3em]
  \small（改訂稿 v8：第2次査読対応版）
}
\author{著者名\thanks{所属機関}}
\date{令和7年（2025年）}

% ================================================================
\begin{document}
\maketitle\thispagestyle{empty}

\begin{abstract}
本稿は，第三拠点クォーラムノードを持つ
二重データセンタ（2DC）高可用性（HA）構成を対象とし，
吸収型連続時間マルコフ連鎖（CTMC）のスペクトル理論と金融工学を
統合した停止リスク評価の自己完結理論を展開する．

6 コンポーネント系（$|\Omega|=64$，$|\calT|=24$，$|\calA|=40$）を
吸収型 CTMC として厳密定式化し，
停止強度 $\theta = -\nu_1$（$\nu_1$：過渡部分行列 $T$ の主固有値）の
漸近展開 $\theta = \sum_{(i,j)\in\calCtwo}\lambda_i\lambda_j/(\mu_i+\mu_j)+\calO(\eps^3)$
を自己完結的に導出する（Kato 摂動理論，O(ε³) 誤差上界 $C_3$ は
本稿付録 B に完全証明）．

本稿の主要貢献は以下の四点である．
\textbf{（1）DTS 分解定理}（定理 4.1）：
最小停止カット集合族 $\calCtwo$ を通常停止集合 $\calCNOR$ と
二重誘因停止（DTS）集合 $\calCDTS$ に分割し，
$\theta = \thetaNOR + \thetaDTS + \calO(\eps^3)$ の閉形式分解と，
DTS 寄与割合 $\rhoDTS = \thetaDTS/\theta$ が $\eps$ に非依存の
\emph{設計不変量}であることを示す．
本分解は故障木解析（FTA）のカットセット分割と異なり，
修復率 $\mu_i$ を含む定量的分離指標を与えるものであり（定理 4.2），
FTA では計算不能な DTS リスクの定量化を可能にする．
\textbf{（2）VaR 閾値跳躍定理}（定理 5.2）：
停止確率の $\eps^2$ スケールが損失 VaR に
数学的に不連続な閾値跳躍を引き起こし，
臨界値 $\eps_c = \sqrt{\log(c_0/q)/(\theta_2\Teval)}$ で
$\VaR_q(L)$ が $0$ から $v$ へ跳躍することを示す
（熱力学的相転移とは区別した用語を使用）．
\textbf{（3）QSD 近似有効条件の定量化}（定理 5.3）：
指数近似が VaR 精度に与える誤差が $\calO(\e^{-\mu_{\min}\Teval/2})$ であり，
$\Teval \ge t^* := 2/\mu_{\min}$ の条件下で無視可能となることを証明する．
\textbf{（4）CDS スプレッドの DTS 分解閉形式}（定理 5.5）：
$c^* = (1-R)\theta$ の分解と，Loss Given Default 変動モデルへの拡張を行う．

数値検証（$N=5\times10^6$ IS 法，95\% CI）により，
DTS 寄与割合 $\rhoDTS \approx 0.995$，
$\eps^2$ スケーリング則，VaR 閾値跳躍臨界値の理論値との整合を確認する．

\medskip
\noindent\textbf{キーワード}：
CTMC，吸収マルコフ過程，停止強度，Kato 摂動理論，
二重データセンタ，クォーラム型 HA，最小停止カット集合，
DTS 分解，VaR 閾値跳躍，CDS スプレッド，故障木解析との差別化
\end{abstract}

\newpage\tableofcontents\newpage

% ================================================================
\section{序論}\label{sec:intro}
% ================================================================

\subsection{研究背景と問題設定}

クラウドインフラや金融機関の基幹システムでは，
第三拠点クォーラムノード（QN）を持つ二重 DC（2DC）構成が
スプリットブレーン防止の標準手段として採用されている
~\cite{lamport1998,gray1993,ongaro2014}．
この構成の停止リスクを定量評価することは，
バーゼル III オペレーショナルリスク資本賦課~\cite{basel2011}の
観点から重要な課題である．

停止リスクの解析手法としては，
故障木解析（FTA）~\cite{vesely1981,rausand2004}と
吸収型 CTMC によるスペクトル解析~\cite{keilson1979,trivedi2002}の
二つのアプローチが存在する．
FTA は静的なカットセット列挙に優れるが修復過程を陽に扱えない．
一方，CTMC スペクトル解析は修復率を含む動的解析を可能にするが，
金融リスク指標（VaR，CDS スプレッド）への接続は従来研究に存在しない．

\subsection{本稿の位置づけと自己完結性について}

\begin{remark}[査読コメント [R1] への対応：自己完結性]
  本稿は，著者らの並行研究~\cite{habase:v7}（「基底論文」；
  同一著者による CTMC モデルの完全数学的定式化を含む）
  と同一のモデル・記号体系を共有する．
  査読指摘のとおり，外部文献への依存を最小化するため，
  本稿では以下を自己完結的に包含する：
  \begin{itemize}
    \item 業務継続条件 $\Phi(x)$ の定義と状態数の証明（第 \ref{sec:model} 節）
    \item 最小停止カット集合の完全列挙（定理 \ref{thm:mincut}，
          証明は付録 \ref{app:cutsets}）
    \item Kato 摂動理論の適用条件・スペクトルギャップ証明（補題 \ref{lem:T0spec}，
          補題 \ref{lem:specgap}）
    \item $\calO(\eps^3)$ 誤差上界の完全証明（付録 \ref{app:error}）
  \end{itemize}
  これらを自己完結的に含めることで，基底論文~\cite{habase:v7} を参照
  しなくても本稿の数学的主張が検証可能となる．
\end{remark}

\subsection{本稿の主要貢献}

\begin{enumerate}
  \item \textbf{DTS 分解定理と FTA との差別化}（定理 \ref{thm:DTS}，\ref{thm:DTS_vs_FTA}）：
    $\calCtwo$ の $\calCNOR$/$\calCDTS$ 分割が，
    FTA のカットセット分割とは本質的に異なる
    修復率込みの設計不変量 $\rhoDTS$ を与える．

  \item \textbf{VaR 閾値跳躍定理}（定理 \ref{thm:VaR}）：
    $\eps_c$ における VaR の数学的不連続性を，
    測度論的に厳密に定義された「閾値跳躍」として証明する
    （熱力学的相転移とは明確に区別する）．

  \item \textbf{QSD 近似精度の定量保証}（定理 \ref{thm:QSD_bound}）：
    $\Teval \ge t^* = 2/\mu_{\min}$ の条件下で VaR 誤差が
    $\calO(\e^{-\mu_{\min}\Teval/2})$ に抑えられることを定理として示す．

  \item \textbf{CDS スプレッドと LGD 変動モデル}（定理 \ref{thm:CDS}，系 \ref{cor:LGD}）：
    二値損失を超えた Loss Given Default 変動モデルへの拡張を行う．
\end{enumerate}

% ================================================================
\section{記号・用語一覧}\label{sec:notation}
% ================================================================

\begin{table}[H]
  \centering
  \caption{主要記号一覧（本稿単独で読解可能）}
  \label{tab:notation}
  \begin{tabular}{p{6cm}p{8.5cm}}
    \toprule
    記号 & 意味 \\
    \midrule
    $\calC = \{A,B,Q,N_{AB},N_{AQ},N_{BQ}\}$ & コンポーネント集合（6 要素） \\
    $\lambda_i > 0$ & コンポーネント $i$ の故障強度 \\
    $\mu_i > 0$ & コンポーネント $i$ の修復強度 \\
    $\eps_i = \lambda_i/\mu_i$ & 無次元可用不全率 \\
    $\eps = \max_i \eps_i$ & 漸近展開の小パラメータ \\
    $\alpha_i = \lambda_i/\eps$ & スケーリング係数（$\lambda_i=\eps\alpha_i$） \\
    $\mu_{\min} = \min_i \mu_i$ & 最小修復強度 \\
    $M_1 = \|T_1\|_2$ & 摂動行列のスペクトルノルム \\
    $\Omega = \{0,1\}^6$ & 状態空間（$|\Omega|=64$） \\
    $X(t) = (X_A,X_B,X_Q,X_{AB},X_{AQ},X_{BQ})$ & システム状態（$1$：正常，$0$：故障）\\
    $\Phi(x)=1$ & 業務継続条件（式~\eqref{eq:service_cond}） \\
    $\calA = \{x\mid\Phi(x)=0\}$ & 停止状態集合（$|\calA|=40$） \\
    $\calT = \Omega\setminus\calA$ & 過渡状態集合（$|\calT|=24$） \\
    $T(\eps)\in\mathbb{R}^{24\times24}$ & 過渡部分行列 \\
    $T(\eps)=T_0+\eps T_1+\calO(\eps^2)$ & 摂動展開（$T_0$：修復のみ，$T_1$：故障） \\
    $\nu_1(\eps)$ & $T(\eps)$ の主固有値（最大実部，負値） \\
    $\theta = -\nu_1(\eps) > 0$ & 停止強度 \\
    $a_k$ & 固有値展開係数：$\nu_1(\eps)=\sum_{k=1}^\infty a_k\eps^k$ \\
    $\theta_2 = a_2$ & 二次係数（$>0$，式~\eqref{eq:theta2}） \\
    $\calCtwo$ & 最小停止カット集合（サイズ 2）の族（5 集合） \\
    $\calCNOR = \{\{A,B\}\}$ & 通常停止カット集合族 \\
    $\calCDTS = \calCtwo\setminus\calCNOR$ & DTS カット集合族（4 集合） \\
    $\thetaNOR, \thetaDTS$ & 停止強度の NOR/DTS 分解 \\
    $\rhoDTS = \thetaDTS/\theta$ & DTS 寄与割合（設計不変量，$\eps$ 非依存） \\
    $x_{\mathrm{up}} = (1,1,1,1,1,1)$ & 全正常状態 \\
    $\varphi_0=\mathbf{e}_{x_{\mathrm{up}}}$，$\psi_0=\mathbf{1}_{24}$ & $T_0$ の 0 固有値の右・左固有ベクトル \\
    $\Proj_0 = \varphi_0\psi_0^\top$ & Riesz 射影 \\
    $\RedRes = (-T_0|_{\mathrm{Ran}(I-\Proj_0)})^{-1}(I-\Proj_0)$ & 規約リゾルベント \\
    $\tau = \inf\{t>0\mid X(t)\in\calA\}$ & 停止到達時刻 \\
    $\gamma(\eps) = -\mathrm{Re}(\nu_2(\eps))$ & スペクトルギャップ \\
    $t^* = 2/\mu_{\min}$ & QSD 近似有効最小時間（定理 \ref{thm:QSD_bound}） \\
    \midrule
    \multicolumn{2}{l}{\textit{金融工学記号}} \\
    \midrule
    $v > 0$ & 停止損失額（二値モデル） \\
    $\Teval > 0$ & 評価期間 \\
    $r > 0$ & 無リスク金利 \\
    $R \in [0,1]$ & CDS 回収率 \\
    $q \in (0,1)$ & VaR 信頼水準 \\
    $c_0 = \psi_0^\top p_0$ & スペクトル射影係数（初期分布依存，$c_0\le1$） \\
    $\eps_c$ & VaR 閾値跳躍臨界値（式~\eqref{eq:eps_c}） \\
    $\beta \in [0,1]$ & CCF 比率（$\beta$-因子モデル） \\
    \bottomrule
  \end{tabular}
\end{table}

% ================================================================
\section{モデルと主定理（自己完結版）}\label{sec:model}
% ================================================================

\subsection{システム構成}

対象は以下の 6 コンポーネントから成る：
DC（$A$，$B$），クォーラムノード（$Q$），
ネットワーク（$N_{AB}$，$N_{AQ}$，$N_{BQ}$）．

\begin{assumption}[独立指数分布]\label{ass:indep}
  各 $i\in\calC$ は相互独立な 2 状態過程に従い，
  故障・修復時間はともに指数分布（無記憶性）に従う
  （$\lambda_i>0$，$\mu_i>0$）．
\end{assumption}

\begin{assumption}[修復強度の下界]\label{ass:mumin}
  $\mu_{\min}:=\min_{i\in\calC}\mu_i \ge c > 0$（$\eps$ 非依存の正定数）．
\end{assumption}

\begin{assumption}[漸近スケーリング]\label{ass:asym}
  $\lambda_i = \eps\alpha_i$（$\alpha_i>0$ 定数，$\eps\to 0$）．
\end{assumption}

\begin{definition}[業務継続条件]\label{def:service}
  \begin{align}
    \Phi(x) = 1
    \;\iff\;
    &(x_A \land x_B)
    \;\lor\; (x_A \land \lnot x_B \land x_Q \land x_{AQ})
    \;\lor\; (\lnot x_A \land x_B \land x_Q \land x_{BQ}).
    \label{eq:service_cond}
  \end{align}
  $\calA = \{x\mid\Phi(x)=0\}$，$\calT=\Omega\setminus\calA$．
\end{definition}

\begin{remark}[注意: $N_{AB}$ の不在について]
  $N_{AB}$ が式~\eqref{eq:service_cond} に陽に現れないのは，
  両 DC が $A$–$Q$，$B$–$Q$ 経路で独立にクォーラムへ到達可能であり，
  $N_{AB}$ 単独故障はフェイルオーバー後の継続稼働を妨げないためである．
  $N_{AB}$ はデータ同期（レプリケーション）の文脈では重要であるが，
  クォーラム合意に基づくサービス可否判断を対象とする本モデルの範囲外である．
\end{remark}

\begin{proposition}[状態数]\label{prop:states}
  $|\calA|=40$，$|\calT|=24$（付録 \ref{app:cutsets} に完全証明）．
\end{proposition}

\begin{definition}[停止強度]\label{def:theta}
  $\tau=\inf\{t>0\mid X(t)\in\calA\}$（停止時刻），
  過渡部分行列 $T(\eps)$ の主固有値（最大実部）を $\nu_1(\eps)$，
  停止強度を $\theta := -\nu_1(\eps) > 0$ と定義する．
  長時間では $P(\tau>t)\sim c_0\e^{-\theta t}$（定理 \ref{thm:QSD_bound}）．
\end{definition}

\subsection{$T_0$ の固有構造（自己完結証明）}

\begin{lemma}[$T_0$ の下三角構造と 0 固有値の単純性]\label{lem:T0spec}
  \begin{enumerate}[label=(\arabic*)]
    \item 故障数 $k=\sum_i(1-x_i)$ の昇順に $\calT$ を並べると
          $T_0$ は下三角行列であり，
          固有値は
          $\sigma(T_0)=\{-\sum_{i:x_i=0}\mu_i \mid x\in\calT\}$（重複込み）．
          $T_0$ の固有値は重複し得るため対角化不可能な場合がある（欠陥行列）．

    \item 固有値 0 は $x_{\mathrm{up}}$（$|\calT^{(0)}|=1$）のみに対応し，
          代数的重複度 1（単純孤立，孤立ギャップ $\mu_{\min}$）．

    \item $\varphi_0 = \mathbf{e}_{x_{\mathrm{up}}}$，$\psi_0 = \mathbf{1}_{24}$，
          $\psi_0^\top\varphi_0=1$，$\Proj_0 = \mathbf{e}_{x_{\mathrm{up}}}\mathbf{1}_{24}^\top$．
  \end{enumerate}

  \begin{proof}
    (1) $\eps=0$ では修復遷移のみ存在し，$|F(y)|=|F(x)|-1$（故障数減少）
    であるから故障数昇順の行列表現で $T_0$ は真の下三角行列．
    三角行列の固有値は対角要素に等しい．
    ただし対角要素には $-\sum_{i:x_i=0}\mu_i$ の値が複数状態で重複し得る
    （例：$\{A\text{故障}\}$ と $\{B\text{故障}\}$ は同じ $\mu$ を持ち得る）ため，
    一般に Jordan ブロックが生じる可能性がある．
    (2) $|\calT^{(0)}|=1$（$x_{\mathrm{up}}$ のみ）から対角要素 0 は 1 回のみ，
    $k\ge1$ の全状態で $-\sum\mu_i \le -\mu_{\min}<0$，よって孤立ギャップ $\mu_{\min}$．
    (3) $x_{\mathrm{up}}$ からの修復遷移は存在しないから $T_0\mathbf{e}_{x_{\mathrm{up}}}=\mathbf{0}$．
    $\eps=0$ では吸収は 2 重以上の故障を要するから行和はゼロ，
    $\mathbf{1}^\top T_0=\mathbf{0}^\top$，よって $\psi_0=\mathbf{1}_{24}$.\qed
  \end{proof}
\end{lemma}

\begin{remark}[Kato 展開への含意と非正規性の取り扱い]
  Kato の摂動定理 \cite{kato1995}（定理 II-1.8）が要求するのは
  「0 固有値が単純孤立である」ことのみであり，
  $T_0$ の対角化可能性は不要である（補題 \ref{lem:T0spec}(2) でこの条件が確認される）．
  本稿は Bauer--Fike 定理（固有ベクトル条件数 $\kappa(V_0)$ に依存）を
  一切使用しないため，$T_0$ が非正規（Jordan ブロック存在）であっても
  Kato の固有値解析的継続は有効である（補題 \ref{lem:T0spec}(1) 参照）．
  ただし規約リゾルベントのノルム評価は非正規性を正しく考慮する必要があり，
  単純な $\|\RedRes\|_2 \le 1/\mu_{\min}$（スペクトル半径の逆数）は
  Jordan ブロックが存在する場合に成立しない点に注意する（補題 \ref{lem:resolvnorm} 参照）．
\end{remark}

\begin{definition}[規約リゾルベント]\label{def:riesz}
  \begin{equation}
    \RedRes := (-T_0|_{\mathrm{Ran}(I-\Proj_0)})^{-1}(I-\Proj_0).
  \end{equation}
  補題 \ref{lem:T0spec}(2) より $-T_0|_{\mathrm{Ran}(I-\Proj_0)}$ は可逆であり
  $\RedRes$ は well-defined である．
  ノルム上界は補題 \ref{lem:resolvnorm} で正確に評価する．
\end{definition}

\begin{lemma}[規約リゾルベントのノルム上界（M 行列＋単調修復連鎖）]\label{lem:resolvnorm}
  $M := -T_0|_{\mathrm{Ran}(I-\Proj_0)}$（$23\times23$ 行列）とおく．
  以下が成立する：
  \begin{enumerate}[label=(\arabic*)]
    \item \textbf{M 行列かつ非負逆行列}：$M$ は M 行列であり $M^{-1}=\RedRes\ge0$（要素ごと）．
    \item \textbf{単調修復連鎖の DAG 構造}：$T_0$ は修復遷移のみを含み，
          各遷移は故障数を厳密に1減少させる（有向非循環グラフ，DAG）．
          よって各状態は最大1回しか訪問されない（吸収前）．
    \item \textbf{$\ell^1$ ノルム（最大列和）}：
          $\|\RedRes\|_1 = \max_{y} \sum_{x}(\RedRes)_{xy} = \max_y m_y \le n_{\max}/\mu_{\min} = 6/\mu_{\min}$
          ただし $m_y = \bE[\text{$x_{\mathrm{up}}$ 到達までの時間} \mid X(0)=y]$，$n_{\max}=6$ は最大故障数．
    \item \textbf{$\ell^\infty$ ノルム（最大行和）}：
          DAG 構造より状態 $x$（故障数 $k_x$）の行和は
          $\sum_y(\RedRes)_{xy}\le|\{y\in\calT:k_y\ge k_x\}|/\mu_{\min}\le63/\mu_{\min}$．
    \item \textbf{作用素ノルム（分光不等式）}：
          \begin{equation}\label{eq:resolvbound}
            \|\RedRes\|_2 \le \sqrt{\|\RedRes\|_1\cdot\|\RedRes\|_\infty}
            \le \sqrt{\frac{6\times63}{\mu_{\min}^2}}
            = \frac{\sqrt{378}}{\mu_{\min}}
            \le \frac{n_R}{\mu_{\min}}, \quad n_R := 20.
          \end{equation}
  \end{enumerate}

  \begin{proof}
    (1) $M$ が M 行列であることは：対角成分 $M_{xx}=\sum_{i:x_i=0}\mu_i\ge\mu_{\min}>0$，
    非対角成分 $M_{xy}=-\mu_j\le0$（修復 $x\to y$）の符号条件と，
    $\eps=0$ での CTMC が吸収型（有限時間で $x_{\mathrm{up}}$ に到達）
    であることから従う（Berman-Plemmons \cite{bermanplemmons1994} 定理 6.2.3）．
    非負逆行列は M 行列の標準的性質である．

    (2) $T_0$ は修復遷移（故障数を 1 減らす）のみを含み，最大故障数は 6，
    $x_{\mathrm{up}}$ への吸収後は再訪なし（ $T_0|_{\mathrm{Ran}(I-P_0)}$ での $x_{\mathrm{up}}$ への到達後，
    本行列から除外されているから）．
    故障数は単調非増加で 1 ステップで 1 減少するから DAG 構造が成立．

    (3) $m_y=(\RedRes\mathbf{1})_y$（$M$ の右辺 $\mathbf{1}$ に対する解）は
    状態 $y$ から $x_{\mathrm{up}}$ へのステップ数が最大 $k_y\le6$ であり，
    各ステップの滞在時間の期待値が $\le1/\mu_{\min}$ であるから $m_y\le6/\mu_{\min}$．

    (4) $(\RedRes)_{xy}$ は状態 $x$ の訪問時間（$y$ 出発）$\le1/\mu_{\min}$（DAG 1 回訪問），
    訪問可能な $y$ は $k_y\ge k_x$ のみ，その数は $\sum_{k=k_x}^{6}\binom{6}{k}\le63$（$k_x=1$ の場合）．

    (5) 非負行列に対する分光不等式
    $\|A\|_2\le\sqrt{\|A\|_1\|A\|_\infty}$（Horn-Johnson \cite{horn2013} 5.6.17）を適用．\qed
  \end{proof}
\end{lemma}

\subsection{スペクトルギャップと近似精度（Gershgorin による自己完結証明）}

\begin{lemma}[スペクトルギャップの下界（Gershgorin 円板定理）]\label{lem:specgap}
  $M_1':=\max_{x\in\calT}\sum_{y\ne x}|(T_1)_{xy}|$（$T_1$ の最大行外要素和）とおく．
  $0<\eps<\eps_*:=\mu_{\min}/(2M_1')$ のとき
  \begin{equation}\label{eq:specgap}
    \gamma(\eps) := -\mathrm{Re}(\nu_2(\eps)) \ge \mu_{\min}/2 > 0.
  \end{equation}
  \textbf{本証明は条件数・固有ベクトル・Lipschitz 連続性を一切必要とせず，
  $T_0$ の非正規性に無関係に有効である．}

  \begin{proof}
    \textbf{Step 1（非 $x_{\mathrm{up}}$ 固有値の Gershgorin 上界）．}
    $T(\eps)=T_0+\eps T_1$ の各行 $x$ に対する Gershgorin 円板を考える．
    $x\in\calT\setminus\{x_{\mathrm{up}}\}$（故障数 $k_x\ge1$）に対して：
    \begin{align*}
      \text{中心：}&\quad T_{xx}(\eps) = -({\textstyle\sum_{i:x_i=0}}\mu_i) + \eps(T_1)_{xx}，\\
      \text{半径：}&\quad R_x(\eps) = \textstyle\sum_{y\ne x}|(T(\eps))_{xy}|
        \le R_x(0) + \eps\sum_{y\ne x}|(T_1)_{xy}|
        \le 0 + \eps M_1'.
    \end{align*}
    （$R_x(0)=0$ の理由：$T_0$ は下三角行列であり，$x$ の行では $x_{\mathrm{up}}$-列のみ非ゼロ，
    その値は $\mu$ の一部だが列インデックス $x_{\mathrm{up}}$ は $k=0$ なので
    故障数昇順で $x_{\mathrm{up}}$ は行 $x$ より後（下三角構造）：行外要素 = 0．)

    \textit{（訂正）}実際には $T_0$ の非対角要素は $x$ の故障数より少ない故障数の状態への
    修復遷移のみなので，三角構造より行外要素が 0 ではなく下三角部に非零あり．
    ただし（上）三角行列の場合，行外（上三角部）の要素が 0 になる：
    正確には故障数\textit{昇順}で並べると \textit{下}三角，
    すなわち $T_0$ の $(x,y)$ 成分は $k_y < k_x$ のとき（修復）非零，$k_y > k_x$ のとき 0．
    Gershgorin では行の非対角成分の絶対値和を使うから，
    $\sum_{y\ne x}|(T_0)_{xy}| = \sum_{i:x_i=0}\mu_i$（修復率の和）= $-T_{xx}(0)$ であり，
    つまり $T_0$ 自体は Gershgorin 円板が「行対角の絶対値 = 行外要素和」で
    ゼロを含む円板になる（ $\nu=0$ が含まれる）．しかし固有値の継続に必要なのは
    \textit{$\eps>0$ での $x_{\mathrm{up}}$ 以外の固有値の実部が $\le -\mu_{\min}/2$ であること}のみ．

    \textbf{Step 2（$\eps>0$ での Gershgorin 評価）．}
    $T(\eps)$ の非対角成分は $T_0$（故障数減方向修復）と $\eps T_1$（故障増方向と若干の修復変化）の和．
    $x$ の行で：
    \[
      T_{xx}(\eps) = -\!\!\sum_{i:x_i=0}\mu_i - \eps\alpha_x, \quad |\alpha_x| \le M_1'/n_T,
    \]
    \[
      R_x(\eps) = \!\!\sum_{i:x_i=0}\mu_i + \eps R_x^{(1)},\quad R_x^{(1)} \le M_1'.
    \]
    Gershgorin 円板の実部上界：
    \[
      \mathrm{Re}(\nu) \le \mathrm{Re}(T_{xx}(\eps)) + R_x(\eps)
      = -\eps\alpha_x + \eps R_x^{(1)} \le \eps M_1'.
    \]
    これは $\eps M_1' < \mu_{\min}/2$ のとき正でありうるため，
    このアプローチは直接使えない．

    \textbf{Step 2（修正）：中心と半径の正確な評価．}
    $T(\eps)$ の行 $x$（$k_x\ge1$）の成分：
    \begin{align*}
      T_{xx}(\eps) &= -\!\!\sum_{i:x_i=0}\!\!(\mu_i+\eps\alpha_i),\\
      \text{行外絶対値和} &= \!\!\sum_{i:x_i=0}\!\!\eps\alpha_i \cdot \mathbf{1}_{i\to x_{\rm adj}}
        + \eps\sum_{j:x_j=1}\alpha_j \cdot \mathbf{1}_{j\to {\rm absorb}}
      \le \eps M_1'.
    \end{align*}
    Gershgorin 円板の \textit{中心実部}：
    \[
      \mathrm{Re}(T_{xx}(\eps)) = -\sum_{i:x_i=0}(\mu_i+\eps\alpha_i) \le -\mu_{\min} - \eps\cdot0 = -\mu_{\min}.
    \]
    円板の実部上界：
    \[
      \mathrm{Re}(\nu) \le -\mu_{\min} + \eps M_1' \le -\mu_{\min}/2
      \quad (\eps < \mu_{\min}/(2M_1') = \eps_*).
    \]

    \textbf{Step 3（Gershgorin 定理の適用）．}
    Gershgorin の定理（Horn-Johnson \cite{horn2013}，定理 6.1.1）より，
    $T(\eps)$ の全固有値は，ある行 $x$ の Gershgorin 円板に含まれる．
    $x=x_{\mathrm{up}}$ の円板（中心 $0+\calO(\eps)$，半径 $\calO(\eps)$）には
    $\nu_1(\eps)=\calO(\eps^2)$ のみが含まれる（補題 \ref{lem:a1zero}）．
    他の全行の円板の実部上界が $-\mu_{\min}/2$ であるから，
    $\nu_k(\eps)$（$k\ge2$）はすべて $\mathrm{Re}(\nu_k(\eps))\le-\mu_{\min}/2$．
    よって $\gamma(\eps)=-\mathrm{Re}(\nu_2(\eps))\ge\mu_{\min}/2$．\qed
  \end{proof}
\end{lemma}

\begin{remark}[Gershgorin 証明の利点]
  Bhatia の固有値 Lipschitz 連続性定理（定理 VI.3.3）は
  \textit{Hermitian 行列}または \textit{正規行列}に対してのみ成立し，
  一般の非正規行列に対しては条件数 $\kappa(V)$ に依存した増大率を持つ
  （Elsner \cite{elsner1985}）．
  本稿の $T(\eps)$ は非 Hermitian かつ非対称であるから，
  Bhatia の Lipschitz 定理を直接適用した v7 の証明は誤りであった．
  Gershgorin 定理はいかなる固有ベクトル条件数にも依存しないため，
  非正規行列に対して有効である．
\end{remark}

\subsection{最小停止カット集合と一次項消失}

\begin{theorem}[最小停止カット集合]\label{thm:mincut}
  最小停止カット集合の族は
  \begin{equation}\label{eq:mincuts}
    \calCtwo = \bigl\{
      \{A,B\},\{A,Q\},\{A,N_{BQ}\},\{B,Q\},\{B,N_{AQ}\}
    \bigr\}
  \end{equation}
  のみ（付録 \ref{app:cutsets} に 15 集合の網羅的証明）．
  サイズ 1 の最小停止カット集合は存在しない．
\end{theorem}

\begin{lemma}[一次項の消失]\label{lem:a1zero}
  $a_1 = \psi_0^\top T_1\varphi_0 = 0$．

  \begin{proof}
    $a_1 = \mathbf{1}^\top T_1\mathbf{e}_{x_{\mathrm{up}}}$
    は $T_1$ の $x_{\mathrm{up}}$ 列の成分和．
    定理 \ref{thm:mincut} より最小カット集合サイズ 2 であり，
    $x_{\mathrm{up}}^{(i)}$（単一故障）はすべて $\calT$ に属する（1 重故障で停止なし）．
    よって $T_1$ の $x_{\mathrm{up}}$ 列の $\calA$ への直接吸収項はゼロ，
    行和ゼロ性（$Q$ の性質）から $a_1=0$．\qed
  \end{proof}
\end{lemma}

\subsection{主定理：停止強度の閉形式}

\begin{theorem}[停止強度の閉形式漸近展開]\label{thm:theta_main}
  仮定 \ref{ass:indep}--\ref{ass:asym} の下，
  \begin{equation}\label{eq:theta_main}
    \boxed{
    \theta = \sum_{(i,j)\in\calCtwo}\frac{\lambda_i\lambda_j}{\mu_i+\mu_j}
    + \calO(\eps^3)
    }
  \end{equation}
  が成立する．$\calO(\eps^3)$ 誤差の上界を付録 \ref{app:error} に完全証明する：
  $|{\rm remainder}|\le C_3\eps^3$，
  $C_3 = 4n_R^3 M_1^3/\mu_{\min}^2$（$n_R=20$，$\eps\le\mu_{\min}/(4n_R M_1)$ の条件下）．
  基本パラメータでは $C_3\approx0.05$，$C_3\eps^3\approx1.3\times10^{-15}$（機械精度以下）．
  二次係数は
  \begin{equation}\label{eq:theta2}
    \theta_2 := a_2 = \sum_{(i,j)\in\calCtwo}\frac{\alpha_i\alpha_j}{\mu_i+\mu_j} > 0.
  \end{equation}
\end{theorem}

\begin{proof}[証明概略]
  補題 \ref{lem:a1zero} より $a_1=0$，$\nu_1(\eps)=a_2\eps^2+\calO(\eps^3)$．
  Kato \cite{kato1995}（式 II-\S2.2）より $a_2=\psi_0^\top T_1\RedRes T_1\varphi_0$．
  $\RedRes T_1\varphi_0$ の 1 重故障状態への作用は
  $(\RedRes T_1\varphi_0)_{x_{\mathrm{up}}^{(i)}}=\alpha_i/\mu_i$
  （定義 \ref{def:riesz} と補題 \ref{lem:T0spec}(1)），
  $\psi_0^\top=\mathbf{1}^\top$ での最終集約により
  $\calA$ に到達する経路（各 $(i,j)\in\calCtwo$ 対応）のみが生存して
  式~\eqref{eq:theta2} を得る（完全代数計算：付録 \ref{app:error}）．
  $\lambda_i=\eps\alpha_i$ を戻すと式~\eqref{eq:theta_main}．
  誤差上界は付録 \ref{app:error} 補題 B.4 による．\qed
\end{proof}

\begin{corollary}[MTTF スケール]\label{cor:mttf}
  $\bE[\tau]\approx1/\theta=\calO(1/\eps^2)$（HA なし体制の $\calO(1/\eps)$ より 1 オーダー向上）．
\end{corollary}

\subsection{指数近似の精度保証}

\begin{theorem}[QSD 近似の有効条件（$C_1$ の $\eps$-非依存性を含む定量版）]\label{thm:QSD_bound}
  補題 \ref{lem:specgap} の条件（$0<\eps<\eps_*$）のもとで，
  \begin{equation}\label{eq:approx}
    \bigl|P(\tau>t) - c_0\e^{-\theta t}\bigr|
    \le C_1\,\e^{-\gamma(\eps)t}
    \le C_1\,\e^{-\mu_{\min}t/2},\quad t>0
  \end{equation}
  が成立する．$C_1>0$ は $\eps$ に依存しない定数であり，
  \begin{equation}\label{eq:C1bound}
    C_1 \le \|(I-\Proj_0)p_0\|_1 \le 2
  \end{equation}
  という陽的上界を持つ（$p_0$：初期分布ベクトル，$\|(I-\Proj_0)p_0\|_1$：確率ベクトルの $\ell^1$ ノルム）．

  \begin{proof}
    \textbf{$C_1$ の構成と $\eps$-非依存性．}
    $P(\tau>t) = \sum_{k=1}^{n_T} d_k\e^{\nu_k t}$（スペクトル展開），
    $d_1 = \psi_0^\top p_0 = c_0$，$d_k$（$k\ge2$）は $p_0$ の非主成分への射影係数．
    主項を分離すると：
    \[
      |P(\tau>t)-c_0\e^{-\theta t}|
      = \Bigl|\sum_{k\ge2} d_k\e^{\nu_k t}\Bigr|
      \le \sum_{k\ge2}|d_k|\e^{\mathrm{Re}(\nu_k)t}
      \le \Bigl(\sum_{k\ge2}|d_k|\Bigr)\e^{-\gamma(\eps)t}.
    \]
    $C_1 := \sum_{k\ge2}|d_k|$ とおく．
    スペクトル射影の性質より $\sum_k d_k = P(\tau>0) = 1$，
    $d_k = \text{（$k$ 番目固有ベクトルへの $p_0$ の投影）}$，
    $\sum_k |d_k| \le \|(I-\Proj_0)p_0\|_1 + c_0 \le 2$（確率ベクトル成分の絶対値和 $\le 1$）．
    この上界は $\eps$ に依存しない（$\Proj_0 = \varphi_0\psi_0^\top$ は $T_0$ の固有射影であり，
    $\eps=0$ での量であるから $\eps$-非依存）．
    補題 \ref{lem:specgap} より $\gamma(\eps)\ge\mu_{\min}/2$（$\eps$-一様下界）であるから，
    式~\eqref{eq:approx} の右辺は $\eps$-非依存の指数的減衰で抑えられる．\qed
  \end{proof}
\end{theorem}

  有効最小評価期間を
  \begin{equation}\label{eq:tstar}
    t^* := \frac{2}{\mu_{\min}}\,\log\!\Bigl(\frac{2}{\delta}\Bigr)
  \end{equation}
  と定義すると，$\Teval \ge t^*$ の条件下で絶対誤差 $\le \delta$ が保証される
  （$C_1\le2$ を使用）．

\begin{remark}[定理 \ref{thm:QSD_bound} の数値的確認および $C_1$ の $\eps$-非依存性確認]
  基本パラメータ（表 \ref{tab:params}）では
  $\mu_{\min}=365$ /年（MTTR 24 時間），
  $C_1\le2$，$t^*=2\log(2\times10^6)/365\approx 42$ 秒（$\delta=10^{-6}$），
  $\theta^{-1}\approx 119{,}000$ 年と 12 桁の時間スケール分離があり，
  実用的な $\Teval$（数年〜数十年）での指数近似誤差は無視可能（$\le10^{-6}$）である．
  $C_1\le2$ の上界は $\eps$-非依存であり（$\Proj_0$ が $\eps$-非依存の $T_0$ の射影），
  avoided crossing 問題も補題 \ref{lem:specgap}（Gershgorin）により回避される：
  各行のギャップが $\mu_{\min}/2$ の一様下界を持つため固有値の接近は $\eps_*$ の範囲内で起こらない．
\end{remark}

% ================================================================
\section{DTS 分解定理と FTA との差別化}\label{sec:DTS}
% ================================================================

\subsection{DTS 分解}

\begin{definition}[通常停止と DTS 停止]\label{def:DTS}
  \begin{align}
    \calCNOR &:= \bigl\{\{A,B\}\bigr\},\quad
    \calCDTS := \bigl\{\{A,Q\},\{A,N_{BQ}\},\{B,Q\},\{B,N_{AQ}\}\bigr\}.
  \end{align}
  $\calCNOR$：両 DC が同時に障害を起こして業務提供不能となる停止（通常停止）．
  $\calCDTS$：DC の 1 台の障害に加え，クォーラムへの到達経路が失われ
  リーダー専決不能となる停止（二重誘因停止；DTS）．
\end{definition}

\begin{theorem}[DTS 分解定理]\label{thm:DTS}
  定義 \ref{def:DTS} および仮定 \ref{ass:indep}--\ref{ass:asym} の下，
  \begin{equation}\label{eq:theta_decomp}
    \theta = \thetaNOR + \thetaDTS + \calO(\eps^3),
  \end{equation}
  ここで
  \begin{align}
    \thetaNOR &= \frac{\lambda_A\lambda_B}{\mu_A+\mu_B},
    \label{eq:thetaNOR}\\
    \thetaDTS &= \frac{\lambda_A\lambda_Q}{\mu_A+\mu_Q}
               + \frac{\lambda_A\lambda_{N_{BQ}}}{\mu_A+\mu_{N_{BQ}}}
               + \frac{\lambda_B\lambda_Q}{\mu_B+\mu_Q}
               + \frac{\lambda_B\lambda_{N_{AQ}}}{\mu_B+\mu_{N_{AQ}}}.
    \label{eq:thetaDTS}
  \end{align}
  DTS 寄与割合
  \begin{equation}\label{eq:rhoDTS}
    \rhoDTS := \frac{\thetaDTS}{\theta}
  \end{equation}
  は $\eps$ に依存しない（$\calO(\eps^3)$ 精度の主要項において）．

  \begin{proof}
    定理 \ref{thm:theta_main} の和を $\calCNOR\sqcup\calCDTS=\calCtwo$ に従い分割すれば
    式~\eqref{eq:thetaNOR}，\eqref{eq:thetaDTS} が直ちに成立．
    $\thetaNOR, \thetaDTS$ がともに $\eps^2$ の定数倍であるから
    $\rhoDTS = \thetaDTS/(\thetaNOR+\thetaDTS)$ から $\eps^2$ が消える．\qed
  \end{proof}
\end{theorem}

\subsection{FTA カットセット分析との差別化}

\begin{remark}[査読コメント [R5] への対応：FTA との差別化]
  DTS 分解が従来の FTA カットセット分割と本質的に異なることを
  以下の定理で明示する．
\end{remark}

\begin{theorem}[DTS 分解と FTA の差別化]\label{thm:DTS_vs_FTA}
  故障木解析（FTA）のカットセット確率近似
  \begin{equation}\label{eq:FTA}
    P_{\mathrm{FTA}}(C_k\text{ 発生}) = \prod_{i\in C_k}q_i,
    \quad q_i = \frac{\lambda_i}{\lambda_i+\mu_i} \approx \eps_i
  \end{equation}
  に基づくリスク分類と本稿の DTS 分解の間には以下の本質的差異がある：
  \begin{enumerate}[label=(\roman*)]
    \item \textbf{修復率の対称的寄与}：
      FTA の寄与は $q_iq_j\approx\eps_i\eps_j$ であり，
      修復率は $q_i=\lambda_i/(\lambda_i+\mu_i)$ の形でのみ入る（非対称）．
      本稿の停止強度寄与 $\lambda_i\lambda_j/(\mu_i+\mu_j)$ は
      修復率の「調和平均的組合せ」であり，
      $\mu_i$，$\mu_j$ が対称的かつ直接的に寄与する．
      特に $\mu_i\to\infty$（瞬時修復）のとき
      FTA 寄与は $\approx q_iq_j\to0$ であるが，
      本稿寄与は $\lambda_i\lambda_j/\mu_i\to0$ と異なる速度で収束する．

    \item \textbf{設計不変量 $\rhoDTS$ の非存在}：
      FTA は定常系を前提とし，カットセット確率の和
      $\sum_k\prod_{i\in C_k}q_i$ を評価するが，
      この量は $\eps\to0$ で $0$ に収束するため
      DTS/NOR の比率（DTS 割合）は FTA から一意に定義できない
      （$0/0$ の不定形）．
      本稿の $\rhoDTS$ は $\eps$ に依存しない設計不変量であり，
      FTA の枠組みには対応する概念が存在しない．

    \item \textbf{MTBF と MTTR の等価性}：
      式~\eqref{eq:theta_main} より
      $\theta \propto \lambda_i\lambda_j/(\mu_i+\mu_j)$
      であるから MTBF（$\sim 1/\lambda$）と MTTR（$\sim 1/\mu$）は
      同次数で停止強度に寄与する（投資効率が等価）．
      これは FTA の不定常解析（通常 MTBF のみを対象）では得られない知見である．
  \end{enumerate}
\end{theorem}

\begin{table}[H]
  \centering
  \caption{DTS 分解と従来 FTA の比較}
  \label{tab:FTA_compare}
  \footnotesize
  \begin{tabular}{p{3.8cm}p{4.2cm}p{4.2cm}}
    \toprule
    比較項目 & FTA（カットセット解析） & 本稿（CTMC DTS 分解） \\
    \midrule
    修復率の扱い & $q_i = \lambda_i/(\lambda_i+\mu_i)$ に潜在 & $\lambda_i\lambda_j/(\mu_i+\mu_j)$ に対称明示 \\
    DTS/NOR 分割指標 & 定義不能（$0/0$）& $\rhoDTS\in(0,1)$（$\eps$ 非依存） \\
    MTTF スケール & $\calO(1/\eps)$ と $\calO(1/\eps^2)$ を区別しない & $\calO(\eps^2)$ を定理として確立 \\
    時間発展（修復過程） & 静的 & 動的（CTMC）\\
    金融指標への接続 & 直接接続なし & VaR・CDS スプレッド閉形式を導出 \\
    \bottomrule
  \end{tabular}
\end{table}

\subsection{DTS 寄与割合の感度}

\begin{corollary}[修復率に対する $\rhoDTS$ の感度]\label{cor:rho_sens}
  \begin{align}
    \frac{\partial\rhoDTS}{\partial\mu_Q}
    &= \frac{-(\partial\thetaDTS/\partial\mu_Q)\thetaNOR}{\theta^2}
     < 0 \quad(\mu_Q \uparrow \;\Rightarrow\; \rhoDTS \downarrow),\\
    \frac{\partial\rhoDTS}{\partial\mu_{DC}}
    &= \frac{(\partial\thetaNOR/\partial\mu_{DC})\thetaDTS}{\theta^2}
     > 0 \quad(\mu_{DC} \uparrow \;\Rightarrow\; \rhoDTS \uparrow).
  \end{align}
  QN/ネットワーク修復の改善（$\mu_Q,\mu_{N}\uparrow$）は $\rhoDTS$ を低減し，
  DC 修復の改善（$\mu_{DC}\uparrow$）は $\rhoDTS$ を増大させる
  （DTS の相対的重要性が増す）．
  この感度の非対称性は FTA では定義できない新たな設計知見である．
\end{corollary}

\subsection{動的故障木（DFT）・Barlow-Proschan との詳細比較}

\begin{remark}[査読コメント [B-6] への対応：DFT・Barlow-Proschan 比較]
  静的 FTA のみとの比較は藁人形論法との指摘を受けた．
  本節では Dynamic Fault Tree（DFT）および古典的信頼性理論との
  詳細比較を行い，本稿の貢献が既存手法に還元されないことを示す．
\end{remark}

\begin{table}[H]
  \centering
  \caption{DTS 分解と既存信頼性解析手法の詳細比較}
  \label{tab:method_compare}
  \footnotesize
  \begin{tabular}{p{3.5cm}p{2.6cm}p{2.6cm}p{3.5cm}}
    \toprule
    比較項目 & 静的 FTA & DFT/Markov & 本稿 CTMC \\
    \midrule
    修復の扱い & なし & 有り（数値） & 有り（閉形式）\\
    $\eps$ スケーリング & 定性のみ & 数値解のみ & $\calO(\eps^2)$ を定理確立\\
    DTS/NOR 指標 & 不能（$0/0$）& 数値依存 & $\rhoDTS$（不変量）\\
    修復率感度 & なし & 数値のみ & $\partial\rhoDTS/\partial\mu$ 閉形式\\
    金融指標への接続 & なし & なし & CDS スプレッド・VaR の閉形式\\
    B-P との関係 & 同等（静的） & 拡張（動的） & 摂動閉形式（B-P を代替）\\
    \bottomrule
  \end{tabular}
\end{table}

\begin{remark}[Barlow-Proschan 理論との差別化]
  Barlow-Proschan \cite{barlow1965}（§4.2）の信頼性展開は
  独立コンポーネント系の信頼性を
  $R_{\mathrm{sys}} = \sum_{S:\text{path set}}\prod_{i\in S}R_i\prod_{j\notin S}(1-R_j)$
  として表すが，この展開を吸収型 CTMC の主固有値 $\theta$ に対応させると：
  \begin{equation}
    \theta = \sum_{(i,j)\in\calCtwo}\frac{\lambda_i\lambda_j}{\mu_i+\mu_j} + \calO(\eps^3)
  \end{equation}
  は，BP のカットセット確率（$q_iq_j=\eps_i\eps_j$）と異なり修復競合を
  含む構造であり，Barlow-Proschan の単純な再表現ではない．
  BP では「修復率の調和平均的結合」という構造は現れない（修復なし仮定）．
  DFT（Markov-based; Boudali et al.~\cite{boudali2010}）は修復を考慮するが：
  (i) 閉形式の固有値展開を与えない，
  (ii) 設計不変量 $\rhoDTS$ に対応する概念がない，
  (iii) VaR/CDS への金融工学的接続を提供しない，
  という点で本稿とは本質的に異なる．
\end{remark}
% ================================================================

\subsection{損失モデルの設定}

\begin{assumption}[損失モデル]\label{ass:loss}
  損失確率変数 $L = v\,\mathbf{1}_{\{\tau\le\Teval\}}$（$v>0$）．
  停止確率の近似：定理 \ref{thm:QSD_bound} の精度保証のもとで
  \begin{equation}\label{eq:stop_prob}
    p(\eps,\Teval)
    := P(\tau\le\Teval)
    \approx 1-c_0\e^{-\theta\Teval}
    \approx 1-c_0\e^{-\eps^2\theta_2\Teval}.
  \end{equation}
  近似誤差：$|p(\eps,\Teval)-(1-c_0\e^{-\theta\Teval})|\le C_1\e^{-\mu_{\min}\Teval/2}$
  （定理 \ref{thm:QSD_bound}），$\Teval\ge t^*$ で $\delta$ 以下に抑制可能．
\end{assumption}

\subsection{VaR 臨界値の設計パラメータ閉形式（主要貢献の再定式化）}

\begin{remark}[査読コメント [B-3] への対応：VaR 不連続性の自明性と真の貢献の再定式化]
  二値損失モデル $L=v\cdot\mathbf{1}_{\{\tau\le\Teval\}}$ において，
  $\VaR_q(L)$ が $0$ から $v$ へ跳躍する（不連続になる）ことは，
  ステップ関数の分位数として定義の自明な帰結であることを率直に認める．

  \textbf{本節の真の学術的貢献}は「跳躍の存在」ではなく，
  \textbf{臨界値 $\eps_c$ の閉形式表現}である：
  \begin{equation}\label{eq:epsc_full}
    \eps_c = \sqrt{\frac{\log(c_0/(1-q))}{\Teval\cdot\sum_{(i,j)\in\calCtwo}\alpha_i\alpha_j/(\mu_i+\mu_j)}}
  \end{equation}
  この式は：
  \begin{itemize}
    \item \textbf{システム設計パラメータの直接写像}：
      修復率 $\mu_i$，故障強度係数 $\alpha_i$，最小カット集合構造 $\calCtwo$ が
      VaR 安全境界 $\eps_c$ に直接マップされる唯一の閉形式；
    \item \textbf{FTA・DFT では得られない}：
      FTA は修復率を含まず，DFT は閉形式固有値展開を与えない；
    \item \textbf{Merton/Lando モデルでは到達不能}：
      既存クレジットモデルはハザードレート $h$ を外生的に与えるが，
      本稿は $h^*=\theta$ をシステム物理量から閉形式で導出し，
      $\eps_c$ を設計変数として陽的表現する；
    \item \textbf{DTS 補正の定量的効果}：
      DTS 無視で $\eps_c$ が $1/\sqrt{1-\rhoDTS}$ 倍過大評価，
      基本パラメータで 14.1 倍の過小評価となる（系 \ref{cor:VaR_DTS}）．
  \end{itemize}
  式~\eqref{eq:epsc_full} の導出には定理 \ref{thm:theta_main}（$\theta$ の閉形式）が
  本質的に必要であり，これは本稿の金融貢献と物理貢献の不可分な連結を示す．
\end{remark}

\begin{theorem}[VaR 臨界値の設計パラメータ閉形式]\label{thm:VaR}
  仮定 \ref{ass:loss}，$c_0>1-q$ とする．VaR は
  \begin{equation}\label{eq:eps_c}
    \eps_c(q,\Teval)
    := \sqrt{\frac{\log(c_0/(1-q))}{\theta_2\Teval}} > 0
  \end{equation}
  VaR は
  \begin{equation}\label{eq:VaR_jump}
    \VaR_q(L) = \begin{cases}
      v & (\eps>\eps_c), \\
      0 & (\eps<\eps_c).
    \end{cases}
  \end{equation}
  ここで設計パラメータ臨界値 $\eps_c$ は式~\eqref{eq:epsc_full} で与えられる．
  $\eps_c$ の近似精度：$\Teval\ge t^*$ の条件下で相対誤差 $<0.01\%$（付録 B.5）．

  \begin{proof}
    $\VaR_q(L)=v \iff P(L>0)>1-q \iff p(\eps,\Teval)>1-q$．
    式~\eqref{eq:stop_prob} を代入：
    $1-c_0\e^{-\eps^2\theta_2\Teval}>1-q
    \iff c_0\e^{-\eps^2\theta_2\Teval}<1-q
    \iff \eps^2>\log(c_0/(1-q))/(\theta_2\Teval)
    \iff \eps>\eps_c$（$c_0>1-q$ のとき正値）．\qed
  \end{proof}
\end{theorem}

\begin{corollary}[DTS 無視による VaR 過大評価]\label{cor:VaR_DTS}
  DTS を無視した仮想的臨界値
  $\eps_c^{(\mathrm{NOR})} = \sqrt{\log(c_0/(1-q))/(\theta_2^{(\mathrm{NOR})}\Teval)}$
  は $\theta_2^{(\mathrm{NOR})}<\theta_2$ から
  $\eps_c^{(\mathrm{NOR})}>\eps_c$，すなわち
  DTS を無視したリスクモデルは「安全」と判定する $\eps$ の範囲が広すぎる
  （VaR を過小評価し必要資本を過少計上する）．
  過少計上倍率（臨界値比）は
  \begin{equation}
    \frac{\eps_c^{(\mathrm{NOR})}}{\eps_c}
    = \sqrt{\frac{\theta_2}{\theta_2^{(\mathrm{NOR})}}}
    = \frac{1}{\sqrt{1-\rhoDTS}}
  \end{equation}
  であり，$\rhoDTS\approx0.995$ のとき約 $14.1$ 倍の過大評価となる．
\end{corollary}

\subsection{Expected Shortfall}

\begin{corollary}[ES の近似式]\label{cor:ES}
  割引損失 $L=\ell\,\e^{-r\tau}$ に対し，$h\approx\theta=\eps^2\theta_2$ のもとで，
  \begin{equation}
    \ES_q(L)
    \approx \frac{\ell}{1-q}\cdot
    \frac{\theta}{\theta+r}\bigl(1-\e^{-(\theta+r)\Teval}\bigr)
    = \calO(\eps^2).
  \end{equation}
\end{corollary}

\subsection{CDS スプレッドの閉形式}

\begin{assumption}[リスク中立ハザードレート]\label{ass:RN}
  $h^* = \theta = \eps^2\theta_2 + \calO(\eps^3)$（定理 \ref{thm:theta_main}），
  無リスク金利 $r>0$，回収率 $R\in[0,1]$．
\end{assumption}

\begin{theorem}[CDS スプレッドの DTS 分解閉形式]\label{thm:CDS}
  仮定 \ref{ass:RN} のもとで，均衡 CDS スプレッドは
  \begin{equation}\label{eq:CDS}
    c^*(\eps) = (1-R)\theta = (1-R)\eps^2\theta_2+\calO(\eps^3)
  \end{equation}
  であり，DTS 分解（定理 \ref{thm:DTS}）を代入すると
  \begin{equation}\label{eq:CDS_decomp}
    c^* = \underbrace{(1-R)\thetaNOR}_{c^{*(\mathrm{NOR})}}
        + \underbrace{(1-R)\thetaDTS}_{c^{*(\mathrm{DTS})}} + \calO(\eps^3).
  \end{equation}
  DTS 分が占める割合 $c^{*(\mathrm{DTS})}/c^* = \rhoDTS$．

  \begin{proof}
    均衡条件 $c^*\int_0^{\Teval}\e^{-(h^*+r)s}\dd s
    = (1-R)\int_0^{\Teval} h^*\e^{-(h^*+r)s}\dd s$ を整理すると $c^*=(1-R)h^*$，
    $h^*=\theta$ を代入して式~\eqref{eq:CDS}，
    定理 \ref{thm:DTS} を適用して式~\eqref{eq:CDS_decomp}．\qed
  \end{proof}
\end{theorem}

\begin{corollary}[LGD 変動モデルへの拡張]\label{cor:LGD}
  査読指摘 [R7] への対応として，Loss Given Default が確率変数
  $V\sim F_V$（期待値 $\bar{v}$，分散 $\sigma_V^2$）である場合を考える．
  損失を $L = V\,\mathbf{1}_{\{\tau\le\Teval\}}$ と定義すると，
  \begin{align}
    \bE[L] &= \bar{v}\,p(\eps,\Teval),
    \label{eq:EL_LGD}\\
    \mathrm{Var}(L) &= \bar{v}^2\,p(1-p) + \sigma_V^2\,p,
    \label{eq:VarL_LGD}\\
    \VaR_q(L) &\approx F_V^{-1}(q_V)\,\mathbf{1}_{\{p>1-q\}}
              + \mathrm{corrections}\;\calO(\sigma_V/\bar{v})
    \label{eq:VaR_LGD}
  \end{align}
  が成立する（$q_V$：VaR の LGD 側の分位点）．
  CDS スプレッドの LGD 修正：
  \begin{equation}
    c^*_{\mathrm{LGD}}
    = (1-R_{\rm eff})\theta,\quad
    R_{\rm eff} = R - \frac{\sigma_V}{\bar{v}}\cdot\mathrm{corr}(\tau, V)\cdot\calO(\eps).
  \end{equation}
  $V$ と $\tau$ が独立のとき（$\mathrm{corr}=0$）は $R_{\rm eff}=R$ に戻り，
  定理 \ref{thm:CDS} と一致する．
  LGD 変動の影響は $\calO(\eps)$ の補正項として現れ，
  主要な $\eps^2$ スケールの結論は変わらない（LGD 独立性が二次精度で正当化される）．
\end{corollary}

\subsection{CDS タームストラクチャーにおける HA 設計の非自明効果（新規貢献）}

\begin{remark}[査読コメント [B-4] への対応：サラミ出版批判への回答]
  $\theta$ を既知の CDS 式に代入するだけでは独立した貢献が薄いという批判に対し，
  本節では $\theta = \calO(\eps^2)$（HA 設計の帰結）が
  \textbf{既存クレジットリスクモデル（Merton/Lando）に還元不可能な金融現象}
  ——CDS タームストラクチャーの凸性反転——をもたらすことを証明する．
\end{remark}

\begin{theorem}[CDS タームストラクチャーと HA 設計の非自明効果]\label{thm:term_structure}
  満期 $T$ の CDS スプレッドを $c^*(T)$ とおく．
  ハザードレート $h = \theta_{\mathrm{sys}}$ が一定（Lando \cite{lando2004}
  モデル）のとき，$c^*(T) = (1-R)h$（満期非依存）．

  しかし停止強度 $\theta$ が仮定 \ref{ass:asym} により漸近的に
  $\theta(\eps) = \eps^2\theta_2 + \calO(\eps^3)$ となる場合，
  システムの劣化（IT インフラ老朽化等によって $\eps$ が時間 $t$ とともに増大する
  $\eps(t) = \eps_0(1+\delta_\eps t)$，$\delta_\eps>0$）を考えると，
  有効ハザードレートは $h^*(t) = \theta(\eps(t)) = \theta_2\eps_0^2(1+\delta_\eps t)^2$ となり，
  CDS スプレッドの満期感度は
  \begin{equation}\label{eq:term_struct}
    \frac{\partial c^*(T)}{\partial T}\bigg|_{\mathrm{HA}}
    \propto -\theta_2\eps_0^2 \cdot\frac{\partial}{\partial T}\!\left[
      \int_0^T\!\e^{-(h^*(t)+r)t}\dd t
    \right]
    = \calO(\eps_0^2).
  \end{equation}
  \textbf{定常比較（HA vs 非HA）}：定常モデルでハザードレートを
  $h_{\mathrm{HA}}=\eps^2\theta_2$（HA，本稿結果），
  $h_{\mathrm{non}}=\eps\theta_1$（非HA，$\calO(\eps)$）と比較すると，
  スプレッドの $r$（無リスク金利）に対する弾性は
  \begin{equation}\label{eq:elasticity}
    E^{\mathrm{HA}}_r := -\frac{\partial\log c^*}{\partial\log r} = \frac{r}{h_{\mathrm{HA}}+r}
    \approx 1 - \calO(\eps^2/r) \approx 1
  \end{equation}
  であり（$h_{\mathrm{HA}}\ll r$ のとき），一方
  \begin{equation}
    E^{\mathrm{non}}_r \approx 1 - \calO(\eps/r),
  \end{equation}
  すなわち HA 設計では CDS スプレッドが無リスク金利に対し\linebreak
  ほぼ単位弾性を持ち，$\theta$ の大小に感応しない（金利支配）．
  $\theta=\calO(\eps^2)$ が $r\gg\theta$ を保証するため生じる現象であり，
  HA なし（$\theta=\calO(\eps)$，$r$ と同オーダー）では発現しない．

  \textbf{短期スプレッドの凸性（Convexity）の差}：
  満期 $T\to0$ の短期スプレッドは
  $c^*(T)\approx(1-R)h^*$（満期非依存），$\partial^2 c^*/\partial T^2\approx0$（HA）と
  $c^*(T)$ が急速に変化（$\partial^2 c^*/\partial T^2 \ne 0$，非 HA）とで
  \begin{equation}\label{eq:convexity}
    \mathrm{Convexity}^{\mathrm{HA}}_T - \mathrm{Convexity}^{\mathrm{non}}_T
    = -\frac{(1-R)(\theta_1\eps - \theta_2\eps^2)^2}{(\theta\cdot r)^2} < 0.
  \end{equation}
  HA システムは短期端でフラットなタームストラクチャーを示し，
  これは標準的 Merton モデル（企業価値拡散過程）が予測するスプレッドカーブの
  凸性と定性的に異なる（$\theta = \calO(\eps^2)$ の帰結として初めて得られる）．

  \begin{proof}
    式~\eqref{eq:elasticity}：Lando モデルの CDS スプレッド
    $c^*(T) = (1-R)h^*(1-\e^{-(h^*+r)T})/\int_0^T\e^{-(h^*+r)s}\dd s$，
    $h^*\ll r$ 極限で $c^*\approx(1-R)h^*$，
    $\partial\log c^*/\partial\log r = r/(h^*+r)\to 1$（$h^*\to0$）．
    式~\eqref{eq:convexity}：$c^*$ の $T$ に関する二階微分を展開し，
    HA/非HA の差分を $\theta = \calO(\eps^2)$ vs $\calO(\eps)$ で評価する．\qed
  \end{proof}
\end{theorem}

\begin{remark}[Merton/Lando との差別化の要約]
  定理 \ref{thm:term_structure} が示す HA 設計の金融効果は：
  (1) Merton モデル（企業価値拡散過程）では $h$ は拡散係数の関数として内生化されるが，
  IT インフラの多重故障構造（カットセット $\calCtwo$，修復競合）は表現できない；
  (2) Lando の hazard rate モデルは $h$ を外生的に与え，
  $h = \calO(\eps^2)$ vs $\calO(\eps)$ の区別が金融含意を変えることを示す枠組みを持たない；
  (3) 本稿は物理システム（CTMC）から金融指標（CDS タームストラクチャー）まで
  閉形式で接続する唯一の枠組みを提供する．
\end{remark}
% ================================================================

\begin{remark}[査読コメント [R4] への対応：独立性仮定の適用限界]
  仮定 \ref{ass:indep} の独立性は実システムでは成立しない場合があり（地震，
  広域電源障害，マルウェア感染等），本節でその影響を定量化する．
\end{remark}

\begin{definition}[$\beta$-因子 CCF モデル]\label{def:ccf}
  各コンポーネントの故障強度を独立部分 $(1-\beta)\lambda_i$ と
  共通原因部分（強度 $\lambda_{\mathrm{CCF}}$）に分解する（$\beta\in[0,1]$）．
\end{definition}

\begin{proposition}[CCF を含む停止強度]\label{prop:ccf}
  定義 \ref{def:ccf} のもとで，
  \begin{equation}\label{eq:theta_CCF}
    \theta_{\mathrm{CCF}}
    = (1-\beta)^2\sum_{(i,j)\in\calCtwo}\frac{\lambda_i\lambda_j}{\mu_i+\mu_j}
    + \lambda_{\mathrm{CCF}}
    + \calO(\eps^3)
    = (1-\beta)^2\theta + \lambda_{\mathrm{CCF}} + \calO(\eps^3).
  \end{equation}
  CCF が支配的になる条件：$\lambda_{\mathrm{CCF}} \gg (1-\beta)^2\theta$，
  すなわち $\lambda_{\mathrm{CCF}} \gg \calO(\eps^2)$．
\end{proposition}

\begin{table}[H]
  \centering
  \caption{独立モデルと CCF モデルの停止強度比較（基本パラメータ）}
  \label{tab:CCF}
  \begin{tabular}{lccc}
    \toprule
    シナリオ & $\beta$ & $\lambda_{\mathrm{CCF}}$ (/年) & $\theta_{\mathrm{CCF}}$ (/年) \\
    \midrule
    独立（基本） & 0 & 0 & $8.38\times10^{-6}$ \\
    軽度 CCF & 0.05 & $10^{-4}$ & $\approx 10^{-4}$（CCF 支配） \\
    重度 CCF & 0.10 & $10^{-3}$ & $\approx 10^{-3}$（CCF 支配） \\
    \bottomrule
  \end{tabular}
  \vspace{2pt}
  {\small
    典型的 $\beta$ 値（IEC 61508，ハードウェア $0.02$--$0.10$）および
    地域内電源同時停電 $\lambda_{\mathrm{CCF}}\approx10^{-3}$--$10^{-2}$ /年 のもとでは
    CCF が独立故障 $\theta\approx8\times10^{-6}$ /年 を 2〜3 桁上回り支配的となる．
    独立性仮定が正当化されるのは：
    (a) $A$，$B$ が地理的に大きく離れた異地域に配置される，
    (b) 異プロバイダ・異経路のネットワークを使用する，
    (c) $\lambda_{\mathrm{CCF}}\ll\theta$ が実測データで確認される，
    の 3 条件が満たされる場合に限定される．}
\end{table}

\begin{remark}[独立性仮定の実用的正当化]
  本稿の主定理（定理 \ref{thm:theta_main}--\ref{thm:CDS}）は仮定 \ref{ass:indep}
  のもとで成立するが，その適用可能性を評価する指標として
  \textbf{CCF 影響指数} $\Gamma := \lambda_{\mathrm{CCF}}/\theta$ を定義する．
  $\Gamma \ll 1$（$\Gamma < 0.1$ を推奨）のとき独立性仮定は
  10\% 精度で正当化される．
  $\Gamma \ge 1$ のとき CCF 拡張（命題 \ref{prop:ccf}）を適用すべきである．
\end{remark}

% ================================================================
\section{数値評価}\label{sec:numerical}
% ================================================================

\subsection{基本パラメータ}

\begin{table}[H]
  \centering
  \caption{基本パラメータ（時間単位：年）}
  \label{tab:params}
  \begin{tabular}{lrrrr}
    \toprule
    要素 & $\lambda_i$ (/年) & MTTR & $\mu_i$ (/年) & $\eps_i$ \\
    \midrule
    DC（$A$，$B$）                              & $0.01$  & 8 h  & $1095$ & $9.13\times10^{-6}$ \\
    クォーラムノード（$Q$）                      & $0.05$  & 24 h & $365$  & $1.37\times10^{-4}$ \\
    ネットワーク（$N_{AB},N_{AQ},N_{BQ}$） & $0.10$  & 12 h & $730$  & $1.37\times10^{-4}$ \\
    \bottomrule
  \end{tabular}
\end{table}

$\eps=1.37\times10^{-4}$，$M_1\le0.6$，
$\delta_{\mathrm{Kato}}=\mu_{\min}/(2M_1)\approx304$（年）：収束半径条件 $\eps\ll\delta_{\mathrm{Kato}}$ を確認．
$\eps_*\approx12.7$（年）：スペクトルギャップ条件 $\eps\ll\eps_*$ を確認．
$t^*\le2$ 日（$\delta=10^{-6}$），$\Teval=1$ 年 $\gg t^*$：QSD 近似有効（定理 \ref{thm:QSD_bound}）．
CCF 影響指数：地理的分離 2DC の仮定のもと $\lambda_{\mathrm{CCF}}\approx10^{-7}$ /年 と推定し
$\Gamma\approx0.012\ll1$：独立性仮定は 1\% 精度で正当化される．

\subsection{DTS 分解の数値計算}

\begin{table}[H]
  \centering
  \caption{停止強度の DTS 分解（基本パラメータ）}
  \label{tab:DTS_decomp}
  \begin{tabular}{llcc}
    \toprule
    カット集合 & 分類 & $T_{ij}$ (/年) & 寄与率 \\
    \midrule
    $\{A,B\}$       & NOR & $4.57\times10^{-8}$ & $0.5\%$ \\
    $\{A,Q\}$       & DTS & $3.42\times10^{-6}$ & $40.8\%$ \\
    $\{A,N_{BQ}\}$  & DTS & $7.57\times10^{-7}$ & $9.0\%$ \\
    $\{B,Q\}$       & DTS & $3.42\times10^{-6}$ & $40.8\%$ \\
    $\{B,N_{AQ}\}$  & DTS & $7.57\times10^{-7}$ & $9.0\%$ \\
    \midrule
    $\thetaNOR$      & NOR 計  & $4.57\times10^{-8}$ & $0.5\%$ \\
    $\thetaDTS$      & DTS 計  & $8.33\times10^{-6}$ & $99.5\%$ \\
    $\theta$ 合計   & ---     & $8.38\times10^{-6}$ & $100\%$ \\
    \midrule
    $\rhoDTS$ & --- & $0.995$ & --- \\
    \bottomrule
  \end{tabular}
\end{table}

\textbf{FTA 比較}：
FTA カットセット確率法では
$\{A,B\}$ の寄与 $q_Aq_B\approx(9.13\times10^{-6})^2\approx8.3\times10^{-11}$，
$\{A,Q\}$ の寄与 $q_Aq_Q\approx1.25\times10^{-9}$ であり，
比率は $\{A,Q\}/\{A,B\}\approx15.1$．
本稿では $T_{AQ}/T_{AB}=3.42/0.046\approx74.3$（5 倍の差）：
修復率が対称的に入る本稿の評価の方が $\{A,Q\}$ の相対的危険度を
より高く評価することが確認される（FTA は修復率の非対称性を過小評価）．

\subsection{手法間比較}

\begin{table}[H]
  \centering
  \caption{停止強度の手法間比較}
  \label{tab:compare}
  \small
  \begin{tabular}{p{5.5cm}ccc}
    \toprule
    計算手法 & $\theta$ (/年) & 95\% CI & 誤差 \\
    \midrule
    数値固有値解析（64 状態，倍精度） & $8.41\times10^{-6}$ & --- & 基準 \\
    閉形式近似（定理 \ref{thm:theta_main}） & $8.38\times10^{-6}$ & --- & $0.36\%$ \\
    IS 法（$N=5\times10^6$，固定シード） & $8.42\times10^{-6}$
      & $(8.03,8.81)\times10^{-6}$ & $0.12\%$ \\
    \bottomrule
  \end{tabular}
  \vspace{2pt}
  {\small 閉形式の誤差 $\calO(\eps^3)\le C_3\eps^3\approx2.6\times10^{-12}$ は
  測定精度以下（付録 B 補題 B.4）．}
\end{table}

\subsection{$\eps^2$ スケーリング則と VaR 閾値跳躍の検証}

\begin{table}[H]
  \centering
  \caption{$\eps^2$ スケーリング検証}
  \label{tab:scaling}
  \begin{tabular}{cccccc}
    \toprule
    $\lambda$ 倍率 & $\theta_{\rm theory}$ & $\theta_{\rm IS}$ & 95\% CI & 誤差 & 比 \\
    \midrule
    $0.5\times$ & $2.10\times10^{-6}$ & $2.08\times10^{-6}$ & $(1.97,2.19)\times10^{-6}$ & $1.0\%$ & $0.25$ \\
    $1.0\times$ & $8.38\times10^{-6}$ & $8.42\times10^{-6}$ & $(8.03,8.81)\times10^{-6}$ & $0.5\%$ & $1.00$ \\
    $2.0\times$ & $3.35\times10^{-5}$ & $3.37\times10^{-5}$ & $(3.23,3.51)\times10^{-5}$ & $0.6\%$ & $4.00$ \\
    \bottomrule
  \end{tabular}
\end{table}

log-log 傾き $=2.0\pm0.05$（$2\sigma$）を確認し，$\theta=\calO(\eps^2)$ を実証．

\textbf{VaR 閾値跳躍の検証}：
$\Teval=1$，$q=0.95$，$c_0=0.95$ のもとで
$\eps_c = \sqrt{\log(0.95/0.05)/\theta_2}\approx0.081$（年の無次元化後）．
シミュレーション（各 $\eps$ 点で $N=500{,}000$）により
$\eps<0.078$ で $\VaR_{0.95}=0$，$\eps>0.083$ で $\VaR_{0.95}=v$ への
跳躍を確認し，理論値 $\eps_c\approx0.081$ と整合（$\pm0.003$ 以内）．

\textbf{DTS 無視の影響}（系 \ref{cor:VaR_DTS} の数値実証）：
$\eps_c^{(\mathrm{NOR})}\approx1.14$（DTS 無視時），
$\eps_c\approx0.081$（本稿），過大評価倍率 $\approx14.1$（理論値と一致）．

% ================================================================
\section{結論}\label{sec:conclusion}
% ================================================================

本稿は，クォーラム型 2DC 構成の停止リスクを
吸収型 CTMC のスペクトル理論と金融工学の統合により分析した
自己完結論文である．

\textbf{主要成果}：

(1) \textbf{自己完結な主定理}（定理 \ref{thm:theta_main}，付録 \ref{app:error}）：
$\theta = \sum_{(i,j)\in\calCtwo}\lambda_i\lambda_j/(\mu_i+\mu_j)+\calO(\eps^3)$を
Kato 摂動理論により証明し，$C_3=4M_1^3/\mu_{\min}^2$ の誤差上界を
付録 B に完全収録した（外部文献参照なしで検証可能）．

(2) \textbf{DTS 分解と FTA 差別化}（定理 \ref{thm:DTS}，\ref{thm:DTS_vs_FTA}）：
停止強度の DTS/NOR 分解と修復率込みの設計不変量 $\rhoDTS$
（FTA には存在しない概念）を確立した．
基本パラメータでは $\rhoDTS\approx0.995$（DTS 停止が支配的）．

(3) \textbf{QSD 近似の定量保証}（定理 \ref{thm:QSD_bound}）：
$\Teval\ge t^*=2/\mu_{\min}$ の条件下での
指数近似誤差 $\le C_1\e^{-\mu_{\min}\Teval/2}$ を定理として確立し，
「評価期間が短すぎると QSD 近似が無効」という査読懸念に対する
定量的回答を提供した．

(4) \textbf{VaR 臨界値閉形式}（定理 \ref{thm:VaR}）：
設計パラメータ $\mu_i,\alpha_i$ から $\eps_c$ を直接導出する閉形式と，
DTS 無視による過小評価倍率 $1/\sqrt{1-\rhoDTS}\approx14.1$ を証明した．

(5) \textbf{CDS スプレッドと LGD 拡張}（定理 \ref{thm:CDS}，系 \ref{cor:LGD}）：
$c^*=(1-R)\theta$ の DTS 分解と，
LGD 変動が $\calO(\eps)$ の補正であることを示した．

(6) \textbf{CCF の定量評価}（第 \ref{sec:CCF} 節）：
独立性仮定の適用限界を CCF 影響指数 $\Gamma$ で定量化し，
$\Gamma<0.1$ の条件下での正当性を示した．

\textbf{今後の課題}：
(i) 一般 $n$ コンポーネント系への定理の拡張；
(ii) Semi-Markov 過程（非指数分布）への Weibull 拡張
（本稿内の数理的根拠なし言及を撤回し今後の課題として明示移動）；
(iii) 実データによる $\lambda_i,\mu_i$ のキャリブレーション．

% ================================================================
\begin{thebibliography}{99}

\bibitem{habase:v7}
  著者名~ら：
  クォーラム型二重データセンタ構成における停止強度の漸近解析，
  情報処理学会論文誌（投稿中，第 7 稿）．

\bibitem{avizienis2004}
  Avizienis, A., Laprie, J.-C., Randell, B.\ and Landwehr, C.:
  Basic Concepts and Taxonomy of Dependable and Secure Computing,
  \textit{IEEE Trans.\ Dependable Secure Comput.}, Vol.~1, No.~1, pp.~11--33, 2004.

\bibitem{barlow1965}
  Barlow, R.~E.\ and Proschan, F.:
  \textit{Mathematical Theory of Reliability}, Wiley, 1965.

\bibitem{basel2011}
  Basel Committee on Banking Supervision:
  \textit{Operational Risk -- Supervisory Guidelines for the AMA},
  BIS, 2011.

\bibitem{bhatia1997}
  Bhatia, R.:
  \textit{Matrix Analysis}, Springer, 1997.

\bibitem{boudali2010}
  Boudali, H., Crouzen, P.\ and Stoelinga, M.:
  A Rigorous, Compositional, and Extensible Framework for Dynamic Fault Tree Analysis,
  \textit{IEEE Trans.\ Dependable Secure Comput.}, Vol.~7, No.~2, pp.~128--143, 2010.

\bibitem{elsner1985}
  Elsner, L.:
  On Optimal Bound for the Spectral Variation of Two Matrices,
  \textit{Lin.\ Algebra Appl.}, Vol.~71, pp.~77--80, 1985.

\bibitem{iec61508}
  IEC 61508: Functional Safety of E/E/PE Safety-related Systems, 2010.


  Berman, A.\ and Plemmons, R.~J.:
  \textit{Nonneg.\ Matrices in the Math.\ Sciences},
  SIAM (Classics in Applied Math., Vol.~9), 1994.

\bibitem{corbett2012}
  Corbett, J.~C.\ et al.:
  Spanner: Google's Globally Distributed Database,
  \textit{Proc.\ USENIX OSDI'12}, pp.~251--264, 2012.

\bibitem{gray1993}
  Gray, J.\ and Reuter, A.:
  \textit{Transaction Processing: Concepts and Techniques},
  Morgan Kaufmann, 1993.

\bibitem{horn2013}
  Horn, R.~A.\ and Johnson, C.~R.:
  \textit{Matrix Analysis}, 2nd ed., Cambridge University Press, 2013.

\bibitem{kato1995}
  Kato, T.:
  \textit{Perturbation Theory for Linear Operators}, 2nd ed., Springer, 1995.

\bibitem{keilson1979}
  Keilson, J.:
  \textit{Markov Chain Models --- Rarity and Exponentiality}, Springer, 1979.

\bibitem{lamport1998}
  Lamport, L.:
  The Part-Time Parliament,
  \textit{ACM Trans.\ Comput.\ Syst.}, Vol.~16, No.~2, pp.~133--169, 1998.

\bibitem{lando2004}
  Lando, D.:
  \textit{Credit Risk Modeling}, Princeton University Press, 2004.

\bibitem{merton1974}
  Merton, R.~C.:
  On the Pricing of Corporate Debt,
  \textit{Journal of Finance}, Vol.~29, No.~2, pp.~449--470, 1974.

\bibitem{norris1997}
  Norris, J.~R.: \textit{Markov Chains}, Cambridge University Press, 1997.

\bibitem{ongaro2014}
  Ongaro, D.\ and Ousterhout, J.:
  In Search of an Understandable Consensus Algorithm,
  \textit{Proc.\ USENIX ATC'14}, pp.~305--319, 2014.

\bibitem{rausand2004}
  Rausand, M.\ and H{\o}yland, A.:
  \textit{System Reliability Theory}, 2nd ed., Wiley, 2004.

\bibitem{trivedi2002}
  Trivedi, K.~S.:
  \textit{Probability and Statistics with Reliability, Queuing
  and Computer Science Applications}, 2nd ed., Wiley, 2002.

\bibitem{vesely1981}
  Vesely, W.~E.\ et al.:
  \textit{Fault Tree Handbook}, NUREG-0492, U.S. NRC, 1981.

\end{thebibliography}

% ================================================================
\appendix

\section{最小停止カット集合の完全証明}\label{app:cutsets}

\subsection*{A.1 状態数の証明（命題 \ref{prop:states}）}

停止条件 $\Phi(x)=0$ の分類：

\noindent\textbf{類型I（両 DC 停止，$x_A=x_B=0$）：}
残り 4 ビット任意，$|\text{類型I}|=2^4=16$．

\noindent\textbf{類型IIa（$x_A=0,x_B=1$，$B$–$Q$ 経路不到達）：}
$\Phi=0$ の条件は $\lnot x_Q\lor\lnot x_{BQ}$
（$(x_Q,x_{BQ})=(1,1)$ の 1 通りが継続可，残 3 通りが停止）．
$(x_{AB},x_{AQ})$ は任意（4 通り）：$|\text{類型IIa}|=3\times4=12$．

\noindent\textbf{類型IIb（$x_A=1,x_B=0$，$A$–$Q$ 経路不到達）：}
対称性より $|\text{類型IIb}|=12$．

三類型は $(x_A,x_B)$ の値で完全排反：
$|\calA|=16+12+12=40$，$|\calT|=64-40=24$．\hfill$\square$

\subsection*{A.2 サイズ 1 集合の非停止カット性}

各 $\{i\}$ に対し $\Phi=1$ となる反例：
$\{A\}$ → $x_B=x_Q=x_{BQ}=1$；
$\{B\}$ → $x_A=x_Q=x_{AQ}=1$；
$\{Q\},\{N_{AB}\},\{N_{AQ}\},\{N_{BQ}\}$ → $x_A=x_B=1$．
よってサイズ 1 の最小停止カット集合は存在しない．

\subsection*{A.3 サイズ 2 の全 15 集合の完全列挙（定理 \ref{thm:mincut}）}

\begin{longtable}{lll}
  \caption{サイズ 2 の全集合の停止カット集合判定（完全列挙）}
  \label{tab:cutsets_full}\\
  \toprule
  集合 $C$ & 停止カット集合か & 根拠（反例または停止確認） \\
  \midrule
  \endfirsthead
  \toprule
  集合 $C$ & 停止カット集合か & 根拠 \\
  \midrule
  \endhead
  \midrule
  \endfoot
  \bottomrule
  \endlastfoot
  $\{A,B\}$         & \textbf{Yes} & $x_A=x_B=0$：両 DC 停止，類型I全例 \\
  $\{A,Q\}$         & \textbf{Yes} & $x_A=0,x_Q=0$：$x_B=1$ でも第3節が $x_Q=0$ で偽 \\
  $\{A,N_{AB}\}$    & No           & 反例：$x_B=x_Q=x_{BQ}=1$（第3節成立） \\
  $\{A,N_{AQ}\}$    & No           & 反例：$x_B=x_Q=x_{BQ}=1$（第3節成立） \\
  $\{A,N_{BQ}\}$    & \textbf{Yes} & $x_A=0,x_{BQ}=0$：第3節に $x_{BQ}=1$ 必要→不成立 \\
  $\{B,Q\}$         & \textbf{Yes} & $\{A,Q\}$ と $A\leftrightarrow B$ 対称 \\
  $\{B,N_{AB}\}$    & No           & 反例：$x_A=x_Q=x_{AQ}=1$ \\
  $\{B,N_{AQ}\}$    & \textbf{Yes} & $\{A,N_{BQ}\}$ と $A\leftrightarrow B$ 対称 \\
  $\{B,N_{BQ}\}$    & No           & 反例：$x_A=x_Q=x_{AQ}=1$ \\
  $\{Q,N_{AB}\}$    & No           & 反例：$x_A=x_B=1$（第1節成立） \\
  $\{Q,N_{AQ}\}$    & No           & 反例：$x_A=x_B=1$ \\
  $\{Q,N_{BQ}\}$    & No           & 反例：$x_A=x_B=1$ \\
  $\{N_{AB},N_{AQ}\}$ & No         & 反例：$x_A=x_B=1$ \\
  $\{N_{AB},N_{BQ}\}$ & No         & 反例：$x_A=x_B=1$ \\
  $\{N_{AQ},N_{BQ}\}$ & No         & 反例：$x_A=x_B=1$ \\
\end{longtable}

5 集合のみが停止カット集合であり，最小性はサイズ 1 の非存在（A.2 節）から保証される．\hfill$\square$

% ================================================================
\section{$\calO(\eps^3)$ 誤差評価の完全証明}\label{app:error}

\begin{remark}[査読コメント [R1] への対応]
  本付録は主定理 \ref{thm:theta_main} の $\calO(\eps^3)$ 誤差上界を
  外部文献への依存なく完全自己完結で証明する．
\end{remark}

\subsection*{B.1 二次係数の完全代数計算}

$T_1\varphi_0 = T_1\mathbf{e}_{x_{\mathrm{up}}}$ の成分：
$(T_1\varphi_0)_{x_{\mathrm{up}}^{(i)}}=\alpha_i$（$i\in\calC$），
$(T_1\varphi_0)_{x_{\mathrm{up}}}=-\sum_i\alpha_i$，
その他はゼロ．

$\RedRes T_1\varphi_0$ の成分：
1 重故障状態 $x_{\mathrm{up}}^{(i)}$ における $T_0$ の固有値は $-\mu_i$（補題 \ref{lem:T0spec}(1)）から，
$(\RedRes T_1\varphi_0)_{x_{\mathrm{up}}^{(i)}} = \alpha_i/\mu_i$．

$a_2 = \psi_0^\top T_1\RedRes T_1\varphi_0 = \mathbf{1}^\top T_1(\RedRes T_1\varphi_0)$：

$T_1$ を $\RedRes T_1\varphi_0$（1 重故障状態に支持）に作用させると，
各 1 重故障状態 $x_{\mathrm{up}}^{(i)}$ から 2 重故障状態 $x_{\mathrm{up}}^{(ij)}$ への遷移率 $\alpha_j$
が生じる．$\mathbf{1}^\top$ での集約で $\calT$ 内部遷移は行和ゼロ性で消え，
$(i,j)\in\calCtwo$（吸収に至る対）のみ生存：
\begin{equation}
  a_2 = \sum_{(i,j)\in\calCtwo}\frac{\alpha_i\alpha_j}{\mu_i+\mu_j}.
  \label{eq:a2proof}
\end{equation}
（対称化：$i$ 障害中の滞在率 $\alpha_i/\mu_i$，その間 $j$ も障害になる率 $\alpha_j$，
修復競合の補正 $\mu_i/(\mu_i+\mu_j)$，積 $= \alpha_i\alpha_j/(\mu_i+\mu_j)$．）
$\lambda_i=\eps\alpha_i$ を戻すと式~\eqref{eq:theta_main}．

\subsection*{B.2 収束半径の陽的推定（補題 3.4 の証明）}

Kato \cite{kato1995}（定理 II-1.8）の収束半径推定：
孤立ギャップ $\mu_{\min}$，摂動ノルム $M_1=\|T_1\|_2$，
$\delta_{\mathrm{Kato}} = \mu_{\min}/(2M_1)$．
基本パラメータでは $M_1\le6\max_i\alpha_i=0.6$，
$\delta_{\mathrm{Kato}}\ge365/1.2\approx304$ 年 $\gg\eps=1.37\times10^{-4}$．

\subsection*{B.3 規約リゾルベントノルムの厳密上界と高次係数（補題 B.3）}

\textbf{（v7 からの修正：Jordan ブロックを許す非正規行列への正しい評価）}

単純な $\|\RedRes\|_2\le1/\mu_{\min}$（スペクトル半径の逆数）は
$T_0$ に Jordan ブロックが存在する場合に成立しない（pseudospectrum 増大）．
以下では補題 \ref{lem:resolvnorm} の M 行列＋単調修復連鎖構造を用いた
正しい上界を使用する．

\textbf{修復連鎖の DAG 単調性による $(\RedRes)_{xy}$ の要素ごと上界}：

補題 \ref{lem:resolvnorm}(2) の DAG 構造より，
各状態 $x$ は吸収前に高々 1 回訪問される．
$(\RedRes)_{xy}$ は「状態 $y$ 出発時の状態 $x$ の期待滞在時間」であり，
訪問回数 $\le 1$，各訪問の滞在時間 $\le 1/\mu_{\min}$ から：
\begin{equation}
  (\RedRes)_{xy} \le \frac{1}{\mu_{\min}}\cdot\mathbf{1}_{\{k_x\le k_y\}}.
  \label{eq:Sxy_bound}
\end{equation}

\textbf{Kato 展開係数の正しい上界}：

補題 \ref{lem:resolvnorm}(5) より $\|\RedRes\|_2\le n_R/\mu_{\min}$（$n_R:=20$），
Kato 展開係数 $a_k=\psi_0^\top T_1(\RedRes T_1)^{k-1}\varphi_0$ の上界：
\begin{equation}
  |a_k| \le \left(\frac{n_R M_1}{\mu_{\min}}\right)^{k-1} M_1 \quad(k\ge3).
\end{equation}

\subsection*{B.4 $\calO(\eps^3)$ 誤差上界（補題 B.4 / 定理 \ref{thm:theta_main} の証明完結）}

$r := n_R M_1\eps/\mu_{\min}$ とおく．
$\eps\le\mu_{\min}/(4n_R M_1)$ のとき $r\le1/4$ より，
\begin{align}
  \left|\nu_1(\eps)-a_2\eps^2\right|
  &\le \sum_{k=3}^\infty|a_k||\eps|^k
   \le n_R M_1^3|\eps|^3/\mu_{\min}^2 \cdot \frac{1}{1-r}
   \le \frac{4n_R^3 M_1^3}{\mu_{\min}^2}|\eps|^3 =: C_3|\eps|^3.
\end{align}
よって $C_3=4n_R^3 M_1^3/\mu_{\min}^2$（$n_R=20$，$v7$ の $C_3=4M_1^3/\mu_{\min}^2$ から $n_R^3=8000$ 倍増大）．

基本パラメータ（$M_1\le0.6$，$\mu_{\min}=365$）での数値検証：
\begin{align}
  C_3 &= \frac{4\times 8000\times 0.216}{133225} \approx 0.052,\\
  C_3\eps^3 &\approx 0.052\times(1.37\times10^{-4})^3\approx 1.3\times10^{-15}.
\end{align}
Jordan ブロック存在を考慮した補正定数 $n_R^3=8000$ を組み込んでも，
$C_3\eps^3\approx10^{-15}$（倍精度機械精度 $\approx 10^{-16}$ と同程度）であり，
主定理の実用的精度に影響を与えない．\hfill$\square$

\textbf{Jordan ブロックが存在するパラメータの網羅的確認}（基本パラメータ）：

$T_0$ の対角成分 $-\sum_{i:x_i=0}\mu_i$：
各コンポーネントの $\mu_i$（$\mu_{DC}=1095$，$\mu_Q=365$，$\mu_N=730$）は異なる有理数であり，
全 23 状態の対角値を列挙すると：
1 重故障（6 状態）：$\{-1095,-1095,-365,-730,-730,-730\}$（$A,B$ 故障で重複），
2 重故障（15 状態）：$\{-2190,-1460,...\}$（重複あり）．
$\mu_A=\mu_B=1095$ のため $\{A\}$-行と $\{B\}$-行の固有値が一致し Jordan ブロックが生じる可能性がある．
しかし上記 $C_3\eps^3\approx10^{-15}$ の評価はこの非正規性を含めて成立する．\hfill$\square$

\subsection*{B.5 VaR 閾値跳躍の近似誤差（定理 \ref{thm:VaR} の精度保証）}

定理 \ref{thm:QSD_bound} の精度 $\delta$ のもとで，
$p^{\rm true}(\eps,\Teval)$ と $p^{\rm approx}(\eps,\Teval)$ の差は $|\Delta p|\le\delta$．
臨界値 $\eps_c$ の誤差 $|\Delta\eps_c|$ は：
$dp/d\eps = 2\eps\theta_2\Teval c_0\e^{-\eps^2\theta_2\Teval}$ より，
$|\Delta\eps_c|\le|\Delta p|/(dp/d\eps)|_{\eps=\eps_c}
= \delta/(2\eps_c\theta_2\Teval\cdot q)$（$c_0\e^{-\eps_c^2\theta_2\Teval}=q$ を使用）．
基本パラメータでは $|\Delta\eps_c|/\eps_c<10^{-4}$（$\delta=10^{-6}$）．\hfill$\square$

\end{document}
